% !PRINTED 367
\weberpage{367}
\S.\,98.

\gesperrt{The Aronhold Seven-Systems.}

Of particular importance are the azygetic systems of seven double tangents, which were first considered by \gesperrt{Aronhold}, and which are therefore called \gesperrt{Aronhold seven-systems}. We also call them \gesperrt{complete seven-systems}, or briefly \gesperrt{complete systems}.

That such systems exist is shown by the compilations of the preceding section. For if

\weq{1}{\begin{aligned}
&x_1 y_1,\ x_2 y_2,\ x_3 y_3,\ x_4 y_4,\ x_5 y_5,\ x_6 y_6\\
&x_1 z_1,\ x_2 z_2,\ x_3 z_3,\ x_4 z_4,\ x_5 z_5,\ x_6 z_6
\end{aligned}}

is an azygetic complex pair, then

\[ x_1,\ x_2,\ x_3,\ x_4,\ x_5,\ y_6,\ z_6 \]

is a complete seven-system (because in the complex $y_6 z_6$ none of the $x$ occurs). It will be shown that no azygetic systems of more than seven elements exist. First of all we have the theorem:

7. \gesperrt{Any six elements of a complete system occur in one and only one complex.}

We first prove the second part of the assertion, that is, we prove that, if

\weq{2}{x_1,\ x_2,\ x_3,\ x_4,\ x_5,\ x_6,\ x_7}

is a complete system, then $x_1, x_2, x_3, x_4, x_5, x_6$ cannot occur in two different complexes. If we suppose this to be possible, then the two complexes must form an azygetic pair like (1). In the syzygetic triple

\[\begin{aligned}
&y_1 z_1,\ y_2 z_2,\ y_3 z_3,\ y_4 z_4,\ y_5 z_5,\ y_6 z_6\\
&y_1 y_2,\ z_1 z_2,\ x_1 x_2\\
&y_1 z_2,\ y_2 z_1
\end{aligned}\]

the double tangents $x_3, x_4, x_5, x_6, x_7$ must occur, and since they occur neither in the first nor in the second of these complexes, they must occur in the third. But since only four pairs still remain in this complex, at
%CONT>>
